\documentclass[11pt]{article}
\usepackage[margin=1in]{geometry}
\usepackage{amsmath,amssymb}
\setlength{\parindent}{0pt}
\setlength{\parskip}{6pt}
\begin{document}
\section*{STAT 412 QZ07 --- Question 8}

\subsection*{Problem}

A manufacturer claims that the average tensile strength of thread A exceeds the average tensile strength of thread B by at least 13 kilograms. Researchers wish to test this claim using a 0.01 level of significance. How large should the samples be if the power of the test is to be 0.9 when the true difference between thread types A and B is 6 kilograms? The population standard deviation for thread A is 6.35 kilograms and the population standard deviation for thread B is 5.59 kilograms.

\subsection*{Solution}

This is a one-sided two-sample mean test with known population standard deviations. The distance between the claimed difference and the true difference used for power is
\[
13-6=7.
\]

For significance level $\alpha=0.01$ and power $0.90$,
\[
z_{0.99}=2.326,\qquad z_{0.90}=1.282.
\]

For equal sample sizes,
\[
n=\frac{(z_{0.99}+z_{0.90})^2(\sigma_A^2+\sigma_B^2)}{(13-6)^2}.
\]

Substitute the values:
\[
n=\frac{(2.326+1.282)^2(6.35^2+5.59^2)}{7^2}=19.01.
\]

Round up:
\[
n=20.
\]

\subsection*{Final Answer}

\fbox{The minimum sample size required is 20 for each sample.}

\end{document}
